\subsection{Evidence beyond the Case-Study Framework}
The external SMOTE-ENN experiment extends the control analysis to an established hybrid cleaning method, rather than another variant of the proposed framework. Its default-rule gains over no intervention largely disappear on average after validation-based threshold selection. This corroborates the relevance of operating-point assessment beyond the case-study intervention within the examined synthetic conditions.

Original deletion introduces an additional distinction. A control matching only the surviving synthetic count can give a different conclusion from one that also matches the retained originals. The retention-matched SVM comparisons retain benefits under dispersed corruption in all three geometries, while boundary-near conditions include adverse outcomes. Thus, the evidence does not support a universal explanation that threshold shifts or sample counts account for all hybrid-intervention effects. It instead shows why original retention, synthetic selection, and decision rules require separate controls.

Removal of corrupted observations or endpoints can coexist with removal of valid minority originals. These measurements should accompany predictive metrics when a hybrid method is interpreted as noise handling. The additional experiment uses one configured external hybrid method and the same nine synthetic settings; it neither verifies WRND nor establishes a general result across modern noise-aware oversamplers or real-world datasets. Its design follows earlier findings and requires broader independent confirmation.
